\section{问题一模型的建立与求解}

\subsection{问题一的描述与分析}

问题一给定静止、无风、$-40\,^{\circ}\mathrm C$的低温环境，目标是描述三层复合防护材料中的瞬态传热，并由人体侧表面温度确定可持续暴露时间。三层材料厚度远小于人体表面尺度，且题目将衣料面积与人体表面积按同一量级处理，因此先忽略面内温差，将防护服等效为沿厚度方向串联的三层等面积平板。设$x=0$位于人体侧，$x=L$位于环境侧，$L=d_1+d_2+d_3$，温度场记为$T_i(x,t)$。

功能层位于第二层。题目给出的放热能力随温度变化，但未给出相变前后导热系数、密度或比热的变化关系，因此各层热物性暂按附件常数处理，仅将功能层放热写入能量方程。这样可以保留附件给出的温度--放热曲线，同时避免引入题目没有提供的相界面位置、潜热和两相热物性参数。

\subsection{三层瞬态导热模型}

\subsubsection{控制方程}

对第$i$层应用能量守恒和Fourier导热定律，有
\begin{equation}
\rho_i c_i\frac{\partial T_i}{\partial t}
=k_i\frac{\partial^2T_i}{\partial x^2}+Q_i(T_i),
\qquad i=1,2,3,
\label{eq:q1-governing}
\end{equation}
其中$\rho_i,c_i,k_i$分别为密度、比热容和导热系数。织物层与隔热层不含内部放热，故
\begin{equation}
Q_1=Q_3=0.
\end{equation}

第二层的实验数据给出单位质量放热能力随温度的变化。记由附件数据插值得到的单位质量放热函数为$q_{\rm pcm}(T)$，则功能层体积热源写为
\begin{equation}
Q_2(T)=\rho_2 q_{\rm pcm}(T).
\label{eq:q1-pcm-source}
\end{equation}
若附件采用放热为负的DSC符号约定，则在数据预处理中统一取放热方向为正。由于题目第四问明确规定各温度下放热能力同比例变化，式\eqref{eq:q1-pcm-source}也便于后续直接写成$Q_2^{(4)}=\eta Q_2$。

\subsubsection{初始条件与层间连续条件}

实验者进入低温环境的瞬间，将人体侧及三层衣料统一初始化为
\begin{equation}
T_i(x,0)=37\,^{\circ}\mathrm C,
\qquad i=1,2,3.
\label{eq:q1-initial}
\end{equation}

忽略层间接触热阻，在两个材料界面$x_1=d_1$和$x_2=d_1+d_2$处满足温度连续与热流连续：
\begin{equation}
\left\{
\begin{aligned}
&T_j(x_j^{-},t)=T_{j+1}(x_j^{+},t),\\
&k_j\frac{\partial T_j}{\partial x}(x_j^{-},t)
=k_{j+1}\frac{\partial T_{j+1}}{\partial x}(x_j^{+},t),
\end{aligned}
\right.
\qquad j=1,2.
\label{eq:q1-interface}
\end{equation}

问题一先按基础模型忽略外表面对流边界层热阻，将外界视为恒温冷源，因此环境侧取
\begin{equation}
T_3(L,t)=-40\,^{\circ}\mathrm C.
\label{eq:q1-outer-boundary}
\end{equation}
人体侧温度不固定为$37\,^{\circ}\mathrm C$；$37\,^{\circ}\mathrm C$仅作为初值。人体侧表面温度定义为
\begin{equation}
T_s(t)=T_1(0,t),
\label{eq:q1-skin-temperature}
\end{equation}
并作为本问的主要输出。人体侧能量边界仍需结合题目允许的生理简化闭合，在数值计算前单独确定；不在此处把未知的$T_s(t)$重新指定为恒温边界。

\subsection{微小时间段分离变量求解}

多层材料的热物性不同，直接对全域采用常规分离变量会受到层间未知温度的限制。参考层合材料瞬态传热的拓展分离变量方法\cite{li2018extended}及其在热防护服--空气--皮肤模型中的应用\cite{li2021protective}，将时间轴划分为$[t_n,t_{n+1}]$。在足够小的$\Delta t=t_{n+1}-t_n$内，将两个界面温度作局部线性近似：
\begin{equation}
\begin{aligned}
T_{12}(t)&\approx T_{12}^{n}+\beta_{12}^{n}(t-t_n),\\
T_{23}(t)&\approx T_{23}^{n}+\beta_{23}^{n}(t-t_n).
\end{aligned}
\label{eq:q1-interface-linear}
\end{equation}
其中$\beta_{12}^{n},\beta_{23}^{n}$由界面热流连续条件确定。

功能层的$q_{\rm pcm}(T)$使方程含温度相关源项。为保持每个微小时间段内的线性解析结构，先按该时间段起点的功能层代表温度$\bar T_2^n$冻结热源：
\begin{equation}
q_{\rm pcm}(T_2)\approx q_{\rm pcm}(\bar T_2^n),
\qquad t\in[t_n,t_{n+1}].
\label{eq:q1-source-freeze}
\end{equation}
该近似只作用于单个微小时间段，随后随温度更新。最终通过减小$\Delta t$检验坚持时间是否收敛；若放热峰附近对时间步敏感，再改用局部一阶线性化。

对第$i$层写成
\begin{equation}
T_i(x,t)=w_i(x,t)+u_i(x,t),
\end{equation}
其中$w_i$承担该时间段内的非齐次边界，$u_i$满足齐次边界。展开为
\begin{equation}
u_i(x,t)=\sum_{m=1}^{\infty}B_{i,m}(t)\phi_{i,m}(x),
\end{equation}
即可将每层PDE化为模态系数的一阶方程
\begin{equation}
\frac{\mathrm dB_{i,m}}{\mathrm dt}
+\alpha_i\lambda_{i,m}^{2}B_{i,m}
=F_{i,m}^{n}(t),
\qquad \alpha_i=\frac{k_i}{\rho_i c_i},
\label{eq:q1-mode-ode}
\end{equation}
其中$F_{i,m}^{n}$包含边界齐次化项以及第二层冻结后的放热源。式\eqref{eq:q1-mode-ode}在每个时间段内可直接积分。利用式\eqref{eq:q1-interface}求得界面温度斜率后，将$t_{n+1}$时刻的温度场作为下一时间段初值，循环推进至人体侧表面温度达到给定阈值。

\begin{algorithm}[!htbp]
\caption{恒温人体边界下的微小时间段分离变量求解流程}
\label{alg:q1-thermal}
\KwIn{三层几何与热物性参数，$T_b=37\,^{\circ}\mathrm C$，$h_b=3.0\,\mathrm{W/(m^2\cdot K)}$，$T_\infty=-40\,^{\circ}\mathrm C$，校正DSC曲线$p(T)$，$\beta_{\rm DSC}$，$\Delta t$，$M$}
\KwOut{$T_i(x,t)$，$T_{\rm in}(t)$及$t_{15},t_{10}$}
计算$A_{\rm DSC}=\int p(T)\,\mathrm dT$、$L_{\rm pcm}=60A_{\rm DSC}/\beta_{\rm DSC}$及$\xi_{\rm DSC}(T)$\;
初始化$T_i(x,0)=37\,^{\circ}\mathrm C$，$\xi^0=0$，$n\gets0$\;
\While{$T_{\rm in}(t_n)>10\,^{\circ}\mathrm C$}{
  计算$\bar T_2^n$与$\xi_{\rm trial}=\xi_{\rm DSC}(\bar T_2^n)$\;
  \eIf{$\xi_{\rm trial}>\xi^n$}{
    $c_{2,\rm eff}^n\gets c_2-L_{\rm pcm}\xi'_{\rm DSC}(\bar T_2^n)$，$\xi^{n+1}\gets\xi_{\rm trial}$\;
  }{
    $c_{2,\rm eff}^n\gets c_2$，$\xi^{n+1}\gets\xi^n$\;
  }
  更新$\alpha_2^n=k_2/(\rho_2c_{2,\rm eff}^n)$及$h_e^n=2.38[T_3(L,t_n)-T_\infty]^{0.25}$\;
  线性近似$T_{12},T_{23}$，三层分别作Robin--Dirichlet、Dirichlet--Dirichlet与Dirichlet--Robin本征展开\;
  由两处段末热流连续条件求$\beta_{12}^n,\beta_{23}^n$，推进温度场至$t_{n+1}$\;
  若$T_{\rm in}$首次越过$15\,^{\circ}\mathrm C$或$10\,^{\circ}\mathrm C$，线性插值得$t_{15}$或$t_{10}$\;
  $n\gets n+1$\;
}
\Return{$T_i(x,t)$, $T_{\rm in}(t)$, $t_{15}$, $t_{10}$}\;
\end{algorithm}


\subsection{坚持时间定义与模型输出}

题目给出人体表面温度低于$15\,^{\circ}\mathrm C$时工作困难，低于$10\,^{\circ}\mathrm C$时存在生命危险。故将可持续正常户外工作的时间定义为首次达到$15\,^{\circ}\mathrm C$的时刻
\begin{equation}
t_{15}=\inf\{t>0:T_s(t)\le15\,^{\circ}\mathrm C\},
\label{eq:q1-t15}
\end{equation}
同时计算
\begin{equation}
t_{10}=\inf\{t>0:T_s(t)\le10\,^{\circ}\mathrm C\}
\label{eq:q1-t10}
\end{equation}
作为危险暴露时间。问题一最终输出完整厚度方向温度场$T(x,t)$、人体侧温度曲线$T_s(t)$以及$t_{15},t_{10}$。在获得数值结果前，不预先填写坚持时间。
